%
%Modelo com instruções para preparação de trabalho para o VI EIC 2020
%
%
%

\documentclass[12pt]{article}
\usepackage[utf8]{inputenc}
\usepackage[brazil]{babel}
\usepackage{epsfig}
\usepackage{mathptmx} %para Times Roman
\usepackage{setspace} \onehalfspacing %para espaçamento 1.5
\usepackage{geometry} %para definir a geometria da página
\geometry{left=3cm,right=2cm, top=7cm, bottom=2cm} %margens da página
\setlength{\parindent}{2cm} %tamanho do parágrafo
\usepackage{background}
\backgroundsetup{contents={\includegraphics[scale=1.4]{logo}},
	scale=1,placement=top,opacity=0.5,position={8.3cm,1cm}, vshift=35}
\pagestyle{empty}%retirando numeração

\begin{document}
	
\begin{center}
	\textsc{\textbf{\Large título do minicurso, oficina e atividades culturais: subtítulo (se houver).}}
\end{center}	

\begin{center}
	\begin{minipage}{.4\linewidth}
\begin{center}
	\textbf{João P. Alves}\\ 
	\small Universidade Federal de Alagoas (UFAL)\\ 
	j.p.alves@ufal.br  
\end{center} 
\end{minipage}
\begin{minipage}{.4\linewidth}
\begin{center}
	\textbf{ Maria C. Silva }\\ 
	\small Universidade Federal de Alagoas (UFAL)\\  
	m.c.silva@ufal.br 
\end{center}
\end{minipage}
\end{center}

\begin{center}	
\textbf{Roberto C. Melo}\\ 
\small Instituto Federal de Alagoas (IFAL)\\  
melo.roberto@ifal.br
\end{center}




As propostas dos minicursos, oficinas e atividades culturais a serem submetidos para o VI Encontro de Iniciação Científica, devem ter no máximo três páginas e devem ser escritos na língua portuguesa. Estas instruções foram preparadas no formato em que o trabalho deve ser apresentado. Os nomes, instituição e e-mail dos autores destas instruções são fictícios e estão aí para mostrar o modelo a ser seguido.

As tabelas e ilustrações devem ser citadas no texto, inseridas o mais próximo possível do trecho a que se referem, e padronizadas conforme as ``Normas de Apresentação Tabular" \ do Instituto Brasileiro de Geografia e Estatística (IBGE). Deve-se indicar a fonte consultada (elemento obrigatório, mesmo que seja produção do próprio autor), de acordo com a ABNT NBR 10520:2002 (para todas e quaisquer citações).

As referências devem seguir a norma ABNT NBR 6023:2018. O arquivo com o resumo deve ser enviado em formato PDF. O prazo final para submissão das propostas é XX de novembro de 2020. 


\vspace{2ex}



\begin{center}
	\textbf{REFERÊNCIAS}
\end{center}


\begin{flushleft}
A AVENTURA de Darwin. \textbf{Diário do Nordeste}, Fortaleza, 7 jun. 2006. Caderno 3, p. 4.
	
\

ASSOCIAÇÃO BRASILEIRA DE NORMAS TÉCNICAS. \textbf{ABNT NBR 6023}: informação e documentação: referências: elaboração. Rio de Janeiro: ABNT, 2018.  p.74.

\

ASSOCIAÇÃO BRASILEIRA DE NORMAS TÉCNICAS. \textbf{ABNT NBR 10520}: informação e  documentação: citações em documentos: apresentação. Rio de Janeiro: ABNT, 2002. p. 7.

\


CHIAVENATO, I. Focalizando o novo negócio. In: CHIAVENATO, I. \textbf{Empreendedorismo}: dando asas ao espírito empreendedor. 4. ed. Barueri, SP: Manole, 2012.

\

CONGRESSO DE INICIAÇÃO CIENTÍFICA DA UFPE, 4., 1996, Recife. \textbf{Anais eletrônicos} [...]. Recife: UFPE, 1996. Disponível em: http://www.propesq.ufpe.br/anais/anais.htm. Acesso em: 21 jan. 1997.

\

GUNCHO, M. R. A educação à distância e a biblioteca universitária. In: SEMINÁRIO DE BIBLIOTECAS UNIVERSITÁRIAS, 10., 1998, Fortaleza. \textbf{Anais} [...]. Fortaleza: Tec Treina, 1998, 1 pen drive.

\

NISTA-PICCOLO, V.; TOLEDO, E. (org.). \textbf{Abordagens pedagógicas do esporte}: modalidades convencionais e não convencionais. Campinas, SP: Papirus, 2018. E-book.

\

RODRIGUES, A. L. A. \textbf{Impacto de um programa de exercícios no local de trabalho}. 2009. Dissertação (Mestrado em Fisiopatologia) - Universidade de São Paulo, São Paulo, 2009.

\


SILVA, M. M. L. \textbf{Crimes da era digital}. Net, Rio de Janeiro, nov. 1998. Seção Ponto de vista. Disponível em: http://www.brazilnet.com.br/contexts/brasil/revistas.htm. Acesso em: 28 nov. 1998.
 

\end{flushleft}



\end{document}
